% Fault101 Fault 2_1 Explorer - Step-by-Step Student Guide
% Generated by guide_tools/latex/latexkit.py from the story script; this file is plain LaTeX - edit it directly.
% Build: pdflatex <this file>   (run it twice so the contents page fills in).  Needs scaguide.sty and figures/ next to it.
\documentclass[11pt,letterpaper]{article}
\usepackage{scaguide}
\guidesetup{
  headerleft={ChipWhisperer Fault101  -  Fault 2\_1 Explorer Guide},
  headerright={Voltage Glitching Explorer},
  footerleft={Companion to: Fault 2\_1 - Introduction to Voltage Glitching Explorer.ipynb},
  title={Fault101 Fault 2\_1 Explorer - Step-by-Step Student Guide},
  subject={Step-by-step guide to the Fault 2\_1 Voltage Glitching Explorer notebook: find a working glitch on real hardware by hand}}

\begin{document}

\guidecover{FAULT101  -  FAULT 2\_1  (EXPLORER)}{Voltage Glitching\\ Explorer}{Step-by-step student guide: find a working glitch on your own kit, by hand}

Fault 2\_1 introduces \textbf{voltage glitching}: briefly pulling a chip's power rail down to make it misbehave. This
guide goes with a \textbf{hands-on notebook} that walks you through finding a working glitch on your own
ChipWhisperer-Husky, one experiment at a time - first by firing single glitches yourself and watching what happens,
then by letting the computer search a range of settings for you. \textbf{This lab needs real hardware; there is no
simulated mode.}

\vspace{10pt}

\begin{guidecallout}{tip}{HOW TO USE THIS GUIDE}
This guide goes \textbf{with} the notebook - it does not replace it. Keep the notebook open next to this PDF. Every
experiment tells you what to run, what to change, and what our own kit showed.
\end{guidecallout}

\guidehthree{\textbf{Which files are which}}

{\small\arrayrulecolor{guidegrid}
\rowcolors{2}{white}{guidezebra}
\begin{longtable}{|R{0.4348}|R{0.4493}|R{0.1159}|}
\hline
\rowcolor{guidenavy}\hdrcell{File (folder: jupyter/courses/fault101/)} & \hdrcell{What it is} & \hdrcell{Open it?} \\ \hline
\endhead
\codesmall{Fault 2\_\bk{}1 - Introduction to Voltage Glitching Explorer.\bk{}ipynb} & The Explorer. Everything in this guide happens here. & \textbf{Yes} \\ \hline
\codesmall{SOLN\_\bk{}Fault 2\_\bk{}1 - .\bk{}.\bk{}.\bk{}ipynb} & Answer key for the original (non-Explorer) lab. & If stuck \\ \hline
\codesmall{firmware/\bk{}mcu/\bk{}simpleserial-glitch/\bk{}simpleserial-glitch.\bk{}c} & (a different folder: \codesmall{chipwhisperer/\bk{}firmware/\bk{}.\bk{}.\bk{}.\bk{}}) The target firmware. The function glitched in this lab is \codesmall{glitch\_loop()}: two nested loops that count from 0 to 2500. & Background \\ \hline
\end{longtable}}

\vspace{4pt}

\guidehthree{\textbf{The four colored boxes you will see everywhere}}

\guidelegend

\vspace{4pt}

\guidesmall{Guide version 1.0 - November 2026. Every number in this guide was measured on one real kit:
ChipWhisperer-Husky, a CW313 board and a CW312T-SAM4S target (Windows 11, Python 3.12, ChipWhisperer 6.0.0, firmware
1.5.0). Your kit will very likely need \textbf{different} width/offset numbers - voltage glitching depends on the
exact cabling and board, more than clock glitching does. The \textbf{method} in each experiment (scan, zoom in,
compare) is what carries over; treat every specific number here as an example, not a target to match. \textbf{Settings
can also drift within a single session on the same kit} - read the note after Experiment 2 before you assume a number
here is permanent.}

\clearpage

\tableofcontents

\clearpage

\guidehone{The big picture}

\textbf{What is voltage glitching?} Digital logic needs its supply voltage to stay within a safe range to work
correctly. If you briefly short the supply rail toward ground - for a few nanoseconds, at exactly the wrong moment -
some transistors can fail to switch, which can flip a bit, skip an instruction, or corrupt a value in a register. The
chip keeps running afterwards; it just computed something wrong for an instant. This is called a \textbf{fault}, and
studying how a short voltage glitch creates one is called \textbf{fault injection}: a standard technique for
evaluating whether a device's security logic (a password check, a signature check, a loop that is only supposed to run
once) can be forced down the wrong path by a hardware fault it was never designed to tolerate.

\textbf{What does the test firmware do?} \code{glitch\_loop()} counts from 0 to 2500 in two nested loops and reports
the final count. With no fault it is always exactly 2500. A fault that skips or corrupts an increment, or exits a loop
early, changes that number. Because the correct answer is always the same, you know immediately whether a glitch did
anything: if the reported count is not 2500, something was disturbed.

\textbf{What is different about the Explorer?} The original lab goes straight to the automatic search. This Explorer
has you fire single glitches by hand first, so you can see what \code{normal}, \code{reset} and \code{success}
actually look like before you ask a loop to find thousands of them for you.

\guidehthree{Your path through the notebook:}

\begin{center}
\begin{tikzpicture}
\node[rounded corners=4pt,fill=guidenavy,text=white,align=center,minimum height=1.15cm,minimum width=3.06cm,inner sep=3pt] (n0) at (1.53,0) {\textbf{1  Fire}\\{\footnotesize a glitch by hand}};
\node[rounded corners=4pt,fill=guidenavy,text=white,align=center,minimum height=1.15cm,minimum width=3.06cm,inner sep=3pt] (n1) at (5.14,0) {\textbf{2 - 3  Scan}\\{\footnotesize width and timing}};
\node[rounded corners=4pt,fill=guidenavy,text=white,align=center,minimum height=1.15cm,minimum width=3.06cm,inner sep=3pt] (n2) at (8.75,0) {\textbf{4  Search}\\{\footnotesize automatically}};
\node[rounded corners=4pt,fill=guidenavy,text=white,align=center,minimum height=1.15cm,minimum width=3.06cm,inner sep=3pt] (n3) at (12.36,0) {\textbf{5 - 6  Map}\\{\footnotesize and strengthen}};
\node[rounded corners=4pt,fill=guideteal,text=white,align=center,minimum height=1.15cm,minimum width=3.06cm,inner sep=3pt] (n4) at (15.97,0) {\textbf{7  Measure}\\{\footnotesize reliability}};
\draw[-{Stealth},guidegrey,thick] (n0.east) -- (n1.west);
\draw[-{Stealth},guidegrey,thick] (n1.east) -- (n2.west);
\draw[-{Stealth},guidegrey,thick] (n2.east) -- (n3.west);
\draw[-{Stealth},guidegrey,thick] (n3.east) -- (n4.west);
\end{tikzpicture}
\end{center}

\vspace{6pt}

{\small\arrayrulecolor{guidegrid}
\rowcolors{2}{white}{guidezebra}
\begin{longtable}{|R{0.1159}|R{0.3768}|R{0.5072}|}
\hline
\rowcolor{guidenavy}\hdrcell{Experiment} & \hdrcell{What you do} & \hdrcell{What you learn} \\ \hline
\endhead
\textbf{1} & Fire single glitches by hand at one setting, then a few neighbours. & The three outcomes (normal / reset / success) and the four glitch knobs. \\ \hline
\textbf{2} & Scan the glitch width at a fixed offset. & The target tolerates short glitches and crashes above a `crash edge' - and a window of successes sits just below it. \\ \hline
\textbf{3} & Scan \codesmall{ext\_offset} (which clock cycle is glitched). & Only specific moments are vulnerable; the rest are harmless or just crash the target. \\ \hline
\textbf{4} & Run the automatic search over a range you choose. & Why a GlitchController beats hand-firing, and how the step size trades speed against luck. \\ \hline
\textbf{5} & Draw a width x offset crash map. & The whole picture at once: a slanted crash edge, with successes in a thin band along it. \\ \hline
\textbf{6} & Compare \codesmall{repeat} and the HP/LP transistors. & Whether making the glitch `stronger' actually helps on this kit. \\ \hline
\textbf{7} & Measure how reliable your best setting is. & A single success does not mean a setting works reliably; you need many shots to know its success rate. \\ \hline
\end{longtable}}

\clearpage

\guidehone{Before you start}

Tick every row. Almost every `nothing happens' problem in this lab traces back to something on this list.

{\small\arrayrulecolor{guidegrid}
\rowcolors{2}{white}{guidezebra}
\begin{longtable}{|R{0.0652}|R{0.5217}|R{0.4130}|}
\hline
\rowcolor{guidenavy}\hdrcell{} & \hdrcell{You need} & \hdrcell{How to check} \\ \hline
\endhead
[  ] & \textbf{Lab 0 - SCA101 Setup} passed (same environment is used here). & ChipWhisperer, numpy and matplotlib import without errors. \\ \hline
[  ] & \textbf{ChipWhisperer-Husky} plugged straight into a USB port (not through a hub). & Lab 0's hardware test ends with \codesmall{hardware OK (disconnected again)}. \\ \hline
[  ] & \textbf{CW312T-SAM4S} target seated on the \textbf{CW313} board, connected to the Husky with the 20-pin ribbon. & See the SCA101 Linux guide, Part 4.2, for pictures of this assembly. \\ \hline
[  ] & \textbf{The Husky \codesmall{Crowbar} SMA cable} connected to the CW313's second SMA socket (the \codesmall{Pos}/measurement cable alone is not enough). & Without it, every glitch in this lab will show as \codesmall{normal} - never \codesmall{reset}, never \codesmall{success}. \\ \hline
[  ] & No other notebook holding the Husky open. & \textbf{File \textgreater{} Shut Down Notebook} in any other open notebook before you start this one. \\ \hline
\end{longtable}}

\vspace{4pt}

\begin{guidecallout}{warn}{WATCH OUT}
\textbf{The single most common reason this lab `does nothing':} the Crowbar cable. The \code{Pos} SMA cable only lets
the Husky *measure* power; a separate cable (\code{Crowbar}, to the CW313's other SMA socket) is what actually shorts
the rail. If every shot in Experiment 1 comes back \code{normal}, check this cable before anything else.
\end{guidecallout}

\begin{guidecallout}{tip}{GOOD TO KNOW}
\textbf{You cannot damage the kit by trying `bad' settings.} A \code{reset} just means the target crashed; the
notebook reboots it automatically. If a cell is stopped halfway (the square Stop button), run \code{recover()} in a
new cell to put the glitch hardware and the target back to a known state.
\end{guidecallout}

\guidehone{Experiment 1 - fire a glitch by hand}

\textbf{Goal:} see the three things a single glitch can do, and learn what the four knobs mean.

{\small\arrayrulecolor{guidegrid}
\rowcolors{2}{white}{guidezebra}
\begin{longtable}{|R{0.1594}|R{0.5217}|R{0.3188}|}
\hline
\rowcolor{guidenavy}\hdrcell{Knob} & \hdrcell{What it controls} & \hdrcell{Husky range} \\ \hline
\endhead
\codesmall{WIDTH} & how \textbf{long} the glitch is (how long the power is pulled down). & 0 to 2296 \\ \hline
\codesmall{OFFSET} & \textbf{where inside a clock cycle} the glitch starts. & 0 to 4592 \\ \hline
\codesmall{EXT\_OFFSET} & \textbf{which clock cycle} after the trigger is glitched (0 = the first one). & 0, 1, 2, ... \\ \hline
\codesmall{REPEAT} & how many glitch pulses fire in a row. & 1 to 8192 \\ \hline
\end{longtable}}

\vspace{4pt}

\textbf{Do this:} run the helper cell once (defines \code{fire\_glitch}, \code{fire\_many} and the others), then run
the \textbf{Experiment 1} cell as it is. Change one number at a time and run it again.

{\small\arrayrulecolor{guidegrid}
\rowcolors{2}{white}{guidezebra}
\begin{longtable}{|R{0.1812}|R{0.1014}|R{0.1014}|R{0.1014}|R{0.5145}|}
\hline
\rowcolor{guidenavy}\hdrcell{Try} & \hdrcell{WIDTH} & \hdrcell{OFFSET} & \hdrcell{EXT\_OFFSET} & \hdrcell{What to expect} \\ \hline
\endhead
A (start) & 1770 & 2670 & 1 & right at the edge: expect some of everything \\ \hline
B: a tiny glitch & 100 & 2670 & 1 & too short to do anything \\ \hline
C: a huge glitch & 2200 & 2670 & 1 & too long to survive \\ \hline
D: the wrong moment & 1770 & 2670 & 30 & right length, wrong clock cycle \\ \hline
E: a different offset & 1770 & 1000 & 1 & right length, wrong part of the cycle \\ \hline
\end{longtable}}

\vspace{4pt}

\begin{guidecallout}{see}{YOU SHOULD SEE}
On our kit, \textbf{A} gave a mix of \code{success} and \code{reset}, with hardly any \code{normal} shots - between
about 1 in 4 and 1 in 3 of 30 shots was a success, and nearly all the rest were resets. \textbf{B, D and E} gave only
\code{normal} (too weak, or the right length aimed at the wrong moment). \textbf{C} crashed the target every single
time. Your exact counts will differ - the pattern (A is the interesting one; B/D/E do nothing; C always crashes) is
what to look for.
\end{guidecallout}

\guidehone{Experiment 2 - find the crash edge}

\textbf{Goal:} find how wide a glitch has to be before it starts to disturb the target, and see why a good setting
hides right next to a bad one.

The cell fires \code{SHOTS} glitches at every width from \code{WIDTH\_START} to \code{WIDTH\_STOP} and draws one
stacked bar per width: grey = \code{normal}, red = \code{reset}, green = \code{success}.

\begin{enumerate}[label=\arabic*.,leftmargin=1.8em,itemsep=2pt,topsep=2pt,parsep=0pt,start=1]
\item Run the cell as it is - a coarse pass over the whole width range. Where do the bars turn red? That is the \textbf{crash edge}.
\item \textbf{Zoom in} on the edge (for example \code{WIDTH\_START = 1700}, \code{WIDTH\_STOP = 1850}, \code{WIDTH\_STEP = 10}, \code{SHOTS = 20}) and run again. Do you see green now?
\item Change \code{OFFSET} (try 2600, then 2750) and find the edge again. Does it move, and which way?
\item Change \code{REPEAT} to 2 and find the edge again.
\end{enumerate}

\begin{guidecallout}{see}{YOU SHOULD SEE}
Mostly grey, with \textbf{two separate red bands} - not one simple edge you cross and stay past. On our kit: grey from
width 0 up to about 1750, red from about \textbf{1800-1850}, back to \textbf{grey again} at 1900-2050, then red again
from about 2100 up. Either band can crash the target; the one near 1800 is the one with a success band hiding just
below it (that is where Experiment 1's setting, width 1770, sits). \textbf{No green anywhere in this coarse pass} (not
even in the second red band, at this offset) - the zoomed run around the *first* edge is where green shows up. If your
own kit's bands land somewhere else, zoom in on whichever red region appears first as you raise the width.
\end{guidecallout}

\begin{guidecallout}{warn}{GOOD TO KNOW}
\textbf{These settings drift, even on the same kit, within a single session.} Working through Fault 2\_2 and 2\_3 on
this exact board, a \code{(width, offset, ext\_\bk{}offset)} that gave a clean, repeatable result went completely
quiet 20-30 minutes later with nothing physically changed - board temperature rising from continuous glitching is the
suspected cause (we measured the Husky's internal temperature climb from about 44°C to 53°C over one long session).
Expect the same here: if a width that worked a while ago stops crashing or stops succeeding, don't assume you did
something wrong - rerun Experiment 2 to find the edge again before trusting an old number.
\end{guidecallout}

\guidehone{Experiment 3 - find the right moment}

\textbf{Goal:} a glitch only matters if it lands on the right instruction. \code{EXT\_OFFSET} slides the glitch later
in time, one clock cycle per step, after the trigger.

\begin{enumerate}[label=\arabic*.,leftmargin=1.8em,itemsep=2pt,topsep=2pt,parsep=0pt,start=1]
\item Run the cell as it is: it fires at \code{EXT\_START} to \code{EXT\_STOP}, with a fixed width and offset.
\item Which \code{EXT\_OFFSET} values give green? Which give red? Which give only grey?
\item Widen the range (\code{EXT\_STOP = 60}, \code{SHOTS = 10}) and run again. Does any later cycle give a success?
\item Change \code{WIDTH} by about 10 either way and repeat. Is the vulnerable cycle still the same one?
\end{enumerate}

\begin{guidecallout}{see}{YOU SHOULD SEE}
On our kit, cycle 0 and cycles 4 and up were grey (harmless). Cycles 1 to 3 were the interesting ones - mostly red,
with \textbf{only cycle 1} turning green. Moving off that one cycle made the glitch fail, so timing matters as much as
width. A wider scan found a few more red cycles further out (around 59 to 60) but no second green one.
\end{guidecallout}

\guidehone{Experiment 4 - let the computer search}

You have now found by hand that a working setting is rare and sits at a precise width, offset and clock cycle. Trying
every combination by hand would take forever, so the lab uses a \code{GlitchController} to walk through a range and
keep score automatically.

{\small\arrayrulecolor{guidegrid}
\rowcolors{2}{white}{guidezebra}
\begin{longtable}{|R{0.3768}|R{0.6232}|}
\hline
\rowcolor{guidenavy}\hdrcell{Variable} & \hdrcell{What it does} \\ \hline
\endhead
\codesmall{MAX\_SUCCESSES} & stop the search once this many successes have been found. \\ \hline
\codesmall{NUM\_TRIES} & how many times each (width, offset, ext\_offset) combination is tried. \\ \hline
\codesmall{WIDTH\_RANGE}, \codesmall{OFFSET\_RANGE}, \codesmall{EXT\_OFFSET\_RANGE} & the box of settings to search. \\ \hline
\codesmall{STEP} & the distance between neighbouring width / offset values tried. \\ \hline
\end{longtable}}

\vspace{4pt}

\begin{enumerate}[label=\arabic*.,leftmargin=1.8em,itemsep=2pt,topsep=2pt,parsep=0pt,start=1]
\item Run the search with the ranges as they are. Note the run time and how many attempts it needed.
\item Make \code{STEP} bigger (25), then smaller (1). What happens to the run time and the number of successes found?
\item Raise \code{MAX\_SUCCESSES} to 5.
\item Move a range \textbf{away} from the good area (for example \code{WIDTH\_\bk{}RANGE = (1000, 1040)}). What does the result look like?
\end{enumerate}

\begin{guidecallout}{see}{YOU SHOULD SEE}
\textbf{As it is:} the search stopped at the first success - a hundred to a few hundred attempts, 4 to 15 seconds on
our kit. \textbf{\code{STEP = 25}:} far fewer settings are tried, so whether it lands on a working one is luck -
sometimes a success in about a second, sometimes none in the whole range. \textbf{\code{STEP = 1}:} nothing is
skipped, so it still finds the first success quickly, but asking for more takes much longer. \textbf{A range far from
the good area:} no successes at all, wasted time.
\end{guidecallout}

\guidehone{Experiment 5 - draw a crash map}

\textbf{Goal:} see the whole picture at once. This cell fires glitches over a grid of widths and offsets and plots one
mark per shot.

\begin{enumerate}[label=\arabic*.,leftmargin=1.8em,itemsep=2pt,topsep=2pt,parsep=0pt,start=1]
\item Run the cell as it is (about a minute). Look at the \textbf{shape} of the red region: where does it start? Does the edge run straight, or lean?
\item Is there any offset where a very wide glitch is still harmless? What does that say about where inside a clock cycle the glitch has to land?
\item Zoom in on the edge with a small window and a small step (this run takes a few minutes).
\item Pick three green marks and add WIDTH + OFFSET for each. What do you notice?
\end{enumerate}

\begin{guidecallout}{see}{YOU SHOULD SEE}
The first pass was grey almost everywhere, with a red patch for wide glitches at higher offsets. The zoomed pass
showed a \textbf{slanted} crash-region edge, with green marks lying in a thin line right along it, just outside the
red. On our kit WIDTH + OFFSET was close to the same total (around 4430-4480) for every green mark: a working glitch
needs to be long enough \textbf{and} early enough in the cycle, in just the right combination.
\end{guidecallout}

\guidehone{Experiment 6 - stronger or weaker?}

The glitch is produced by switching a transistor that shorts the target's power rail toward ground. You can try to
make it stronger by repeating the pulse (\code{REPEAT}) or by using both the high-power and low-power transistors
together. A stronger glitch should crash the target at a smaller width - does the crash edge actually move?

{\small\arrayrulecolor{guidegrid}
\rowcolors{2}{white}{guidezebra}
\begin{longtable}{|R{0.1884}|R{0.8116}|}
\hline
\rowcolor{guidenavy}\hdrcell{Setting} & \hdrcell{What it means} \\ \hline
\endhead
\codesmall{hp} & the high-power crowbar transistor. \\ \hline
\codesmall{lp} & the low-power crowbar transistor. \\ \hline
both off & no crowbar at all - a control: nothing should happen. \\ \hline
\end{longtable}}

\vspace{4pt}

\begin{enumerate}[label=\arabic*.,leftmargin=1.8em,itemsep=2pt,topsep=2pt,parsep=0pt,start=1]
\item Run the cell as it is. For every setting in \code{SETTINGS} it scans the width and reports the crash edge.
\item Compare the edges. Does \code{REPEAT} 2 or 4 move it? Does \code{hp + lp}? What does the control do?
\item Add a setting of your own to \code{SETTINGS}.
\end{enumerate}

\begin{guidecallout}{see}{YOU SHOULD SEE}
On our kit, \code{hp} alone at repeat 1, 2 and 4 gave about the \textbf{same} crash edge - repeating the pulse did not
make it stronger here. \code{hp + lp} together moved the edge to a \textbf{smaller} width: two transistors really were
stronger. \code{lp} alone never crashed the target in this range, and the control (both off) never did either. A
surprise like `repeat did nothing' is a result worth writing down, not a mistake.
\end{guidecallout}

\guidehone{Experiment 7 - how reliable is your best glitch?}

\textbf{Goal:} a setting that worked once may not work again. A real evaluation needs to know \textbf{how often} a
given setting succeeds, because that decides how many attempts (and how many crashes) using it actually costs.

\begin{enumerate}[label=\arabic*.,leftmargin=1.8em,itemsep=2pt,topsep=2pt,parsep=0pt,start=1]
\item Run the cell as it is. Which of the listed candidates has the best success rate?
\item Replace the candidates with your own best settings from Experiments 2 to 5.
\item Raise \code{SHOTS} to 200 for the winner and run it twice. How much does the measured rate change by chance alone?
\end{enumerate}

\begin{guidecallout}{see}{YOU SHOULD SEE}
Success rates between 0\% and about 40\% on our kit. Neighbouring settings differed a lot - moving the width or offset
by only 10 steps could take a setting from `works a third of the time' to `never works, always crashes'. The same
setting also gave a slightly different rate each time it was measured, so a single run of 20 shots tells you very
little on its own.
\end{guidecallout}

\guidehone{Challenges (optional)}

Use the helper functions (\code{fire\_glitch}, \code{fire\_many}, \code{scan}, \code{plot\_tallies},
\code{crash\_edge}, \code{glitch\_on}, \code{recover}) and the EDIT ME cells to try these on your own kit:

\begin{enumerate}[label=\arabic*.,leftmargin=1.8em,itemsep=2pt,topsep=2pt,parsep=0pt,start=1]
\item \textbf{The cheapest success.} Find ranges and a \code{STEP} for the automatic search that collect 10 successes in the fewest attempts and the least time.
\item \textbf{Why always the same number?} Look at the counter value after a success. Is it always the same one? What does that tell you about *what* the glitch changed, if the answer is always identical?
\item \textbf{A different vulnerable moment.} Experiment 3 found one vulnerable clock cycle. With a wider scan and different widths, can you find a second one - or one that only ever crashes the target, never succeeds?
\item \textbf{Fewer crashes per success.} Every \code{reset} is a crashed target, which in a real evaluation is easy to notice. Find a setting with a success rate of at least 5\% and fewer than one reset per success.
\end{enumerate}

\guidehone{If something goes wrong}

{\small\arrayrulecolor{guidegrid}
\rowcolors{2}{white}{guidezebra}
\begin{longtable}{|R{0.2754}|R{0.3043}|R{0.4203}|}
\hline
\rowcolor{guidenavy}\hdrcell{What you see} & \hdrcell{Most likely cause} & \hdrcell{What to do} \\ \hline
\endhead
Every shot is \codesmall{normal}, never \codesmall{reset} or \codesmall{success}, even at a large width. & The Crowbar SMA cable is not connected. & Connect it to the CW313's second SMA socket (see Before You Start), then run the setup cells again. \\ \hline
Every shot is \codesmall{reset}. & The width is too large for this offset, or the target is stuck from an earlier crash. & Lower \codesmall{WIDTH} (try 1700), run \codesmall{recover()}, try again. \\ \hline
Plenty of \codesmall{reset}, but no \codesmall{success}. & The settings are slightly off - successes live in a band only a few steps wide, at the edge of the crash region. & Go back to Experiment 2 and zoom in with a step of 5 or 10; use more \codesmall{SHOTS}. \\ \hline
A \codesmall{success} appears once and not the next time at the same setting. & Normal: the success rate is a probability, not a guarantee (see Experiment 7). & Fire more shots and compare *rates*, not single results. \\ \hline
You pressed Stop, or a cell was interrupted. & The target may be mid-crash and the glitch hardware left enabled. & Run \codesmall{recover()} in a new cell. \\ \hline
\codesmall{OSError: Could not find ChipWhisperer} / \codesmall{Unable to communicate with device}. & The Husky is unplugged, or another notebook still holds it. & Plug it in; close every other open notebook (File \textgreater{} Shut Down Notebook); restart this kernel. \\ \hline
\codesmall{USBErrorPipe}, \codesmall{Unknown hwInfoVer}, or \codesmall{ValueError: Unknown ChipWhisperer}. & The Husky's USB link or FPGA stopped answering. & Unplug the Husky's USB cable and plug it back in directly (not through a hub); restart the kernel and run the setup cells again. \\ \hline
A scan or search takes a very long time. & Every \codesmall{reset} costs about half a second while the target reboots. & Use fewer widths (bigger step), fewer \codesmall{SHOTS}, or a narrower range. \\ \hline
\end{longtable}}

\vspace{4pt}

\begin{guidecallout}{warn}{WATCH OUT}
Stay inside the ranges quoted in each experiment, and do not unplug any cable while a cell is running. A crashed
target is a normal outcome here - the notebook reboots it for you automatically.
\end{guidecallout}

\guidehone{Appendix - every variable you can edit}

{\small\arrayrulecolor{guidegrid}
\rowcolors{2}{white}{guidezebra}
\begin{longtable}{|R{0.2174}|R{0.1449}|R{0.6377}|}
\hline
\rowcolor{guidenavy}\hdrcell{Variable} & \hdrcell{Default} & \hdrcell{What it does} \\ \hline
\endhead
\codesmall{WIDTH} & 1770 & Glitch length, in Husky phase-shift steps (0 to 2296). \\ \hline
\codesmall{OFFSET} & 2670 & Where inside the clock cycle the glitch starts (0 to 4592). \\ \hline
\codesmall{EXT\_OFFSET} & 1 & Which clock cycle after the trigger is glitched (0, 1, 2, ...). \\ \hline
\codesmall{REPEAT} & 1 & Number of glitch pulses fired per shot. \\ \hline
\codesmall{SHOTS} & varies & How many glitches to fire for one measurement; more shots = a better estimate of the success rate, but slower. \\ \hline
\codesmall{WIDTH\_\bk{}START/\bk{}STOP/\bk{}STEP} & varies & The width range scanned in Experiments 2, 5 and 6. \\ \hline
\codesmall{OFFSET\_\bk{}START/\bk{}STOP/\bk{}STEP} & varies & The offset range scanned in Experiment 5. \\ \hline
\codesmall{EXT\_START/STOP} & 0 / 8 & The ext\_offset range scanned in Experiment 3. \\ \hline
\codesmall{MAX\_SUCCESSES} & 1 & How many successes the automatic search (Experiment 4) collects before stopping. \\ \hline
\codesmall{NUM\_TRIES} & 5 & How many times each setting is retried in the automatic search. \\ \hline
\codesmall{SETTINGS} & list & (name, repeat, hp, lp) tuples compared in Experiment 6. \\ \hline
\codesmall{CANDIDATES} & list & (width, offset, ext\_offset, repeat) tuples compared in Experiment 7. \\ \hline
\end{longtable}}

\vspace{4pt}

\guidehone{Appendix - Glossary}

{\small\arrayrulecolor{guidegrid}
\rowcolors{2}{white}{guidezebra}
\begin{longtable}{|R{0.2174}|R{0.7826}|}
\hline
\rowcolor{guidenavy}\hdrcell{Term} & \hdrcell{Meaning} \\ \hline
\endhead
Voltage glitch & A brief, deliberate dip in a chip's supply voltage, used to see whether it disturbs normal execution. \\ \hline
Fault injection & Deliberately causing a hardware fault (here, with a voltage glitch) to test how a device's logic responds. \\ \hline
\codesmall{normal} & Outcome: the glitch had no visible effect; the firmware's counter still reached 2500. \\ \hline
\codesmall{reset} & Outcome: the glitch crashed the target; the notebook reboots it and continues. \\ \hline
\codesmall{success} & Outcome: the glitch changed the result (the counter is not 2500). \\ \hline
Crash edge & The boundary (in width, or in width+offset) past which the target reliably crashes instead of running normally. \\ \hline
\codesmall{ext\_offset} & Which clock cycle, counted from the trigger, the glitch is placed on. \\ \hline
Crowbar & The transistor (and its SMA cable on the Husky) that shorts the target's power rail to produce the glitch. \\ \hline
HP / LP & The high-power and low-power crowbar transistors; using both together makes a stronger glitch than either alone. \\ \hline
\codesmall{GlitchController} & The ChipWhisperer class that walks through a grid of glitch settings and tallies the outcomes. \\ \hline
\end{longtable}}

\vspace{4pt}

\guidehthree{Notice}

\guidesmall{This guide accompanies ChipWhisperer course material that references (commercial) training material on
ChipWhisperer.io. You can freely execute and use the lab per the open-source license (including in your own courses if
you distribute similarly), but you must maintain notice about the source: https://www.ChipWhisperer.io.}

\guidesmall{Copyright (C) NewAE Technology Inc., 2015-2020. ChipWhisperer is a trademark of NewAE Technology Inc.,
claimed in all jurisdictions, and registered in at least the United States of America, European Union, and Peoples
Republic of China. Tutorials derived from our open-source work must be released under the associated open-source
license, and notice of the source must be clearly displayed.}

\guidesmall{THE SOFTWARE IS PROVIDED ``AS IS'', WITHOUT WARRANTY OF ANY KIND, EXPRESS OR IMPLIED, INCLUDING BUT NOT
LIMITED TO THE WARRANTIES OF MERCHANTABILITY, FITNESS FOR A PARTICULAR PURPOSE AND NONINFRINGEMENT. IN NO EVENT SHALL
THE AUTHORS OR COPYRIGHT HOLDERS BE LIABLE FOR ANY CLAIM, DAMAGES OR OTHER LIABILITY, WHETHER IN AN ACTION OF
CONTRACT, TORT OR OTHERWISE, ARISING FROM, OUT OF OR IN CONNECTION WITH THE SOFTWARE OR THE USE OR OTHER DEALINGS IN
THE SOFTWARE.}

\end{document}
